On considère la fonction $f$ définie sur l'intervalle $]0~;~+\infty[$ par :

\[f(x) = \dfrac{\ln (x)}{x^{2}} + 1\]

On note $\mathcal{C}_{f}$ la courbe représentative de la fonction $f$ dans un repère orthonormé. On admet que la fonction $f$ est dérivable sur l'intervalle $]0~;~ +\infty[$ et on note $f'$ sa fonction dérivée.

\medskip

\begin{enumerate}
\item Déterminer les limites de la fonction $f$ en $0$ et en $+\infty$.

En déduire les éventuelles asymptotes à la courbe $\mathcal{C}_{f}$.
\item Montrer que, pour tout réel $x$ de l'intervalle $]0~;~ +\infty[$, on a :

\[f'(x) = \dfrac{1 - 2\ln(x)}{x^{3}}\]

\item En déduire le tableau de variation de la fonction $f$ sur l'intervalle $]0~;~ +\infty[$.

\item 
	\begin{enumerate}
		\item Montrer que l'équation $f(x) = 0$ possède une unique solution, notée $\alpha$, sur l'intervalle $]0~;~ +\infty[$.
		\item Donner un encadrement du réel $\alpha$ d'amplitude $0,01$.
		\item En déduire le signe de la fonction $f$ sur l'intervalle $]0~;~ +\infty[$.
	\end{enumerate}
\item On considère la fonction $g$ définie sur l'intervalle $]0~;~ +\infty[$ par :

\[g(x) = \ln(x)\]

On note $\mathcal{C}_{g}$ la courbe représentative de la fonction $g$ dans un repère orthonormé d'origine O. On considère un réel $x$ strictement positif et le point M de la courbe $C_{g}$ d'abscisse $x$. On note OM la distance entre les points O et M.

	\begin{enumerate}
		\item Exprimer la quantité OM$^{2}$ en fonction du réel $x$.
		\item Montrer que, lorsque le réel $x$ parcourt l'intervalle $]0~;~ +\infty[$, la quantité OM$^{2}$ admet un minimum en $\alpha$.
		\item La valeur minimale de la distance OM, lorsque le réel $x$ parcourt l'intervalle $]0~;~+\infty[$, est appelée distance du point O à la courbe $\mathcal{C}_{g}$. On note $d$ cette distance.

Exprimer $d$ à l'aide de $\alpha$.
	\end{enumerate}
\end{enumerate}
